\documentclass[12pt,a4paper]{article}
\usepackage[utf8]{inputenc}
\usepackage[T1]{fontenc}
\usepackage{amsmath}
\usepackage{amsfonts}
\usepackage{amssymb}
\usepackage{graphicx}
\usepackage{geometry}
\usepackage{array}
\usepackage{longtable}
\usepackage{booktabs}
\usepackage{multirow}
\usepackage{hyperref}
\usepackage{xcolor}
\usepackage{tikz}
\usetikzlibrary{shapes,arrows,positioning}

\geometry{left=2.5cm,right=2.5cm,top=2.5cm,bottom=2.5cm}

\title{DetectifAI: AI-Powered CCTV Footage Analyzer\\
Mid-Evaluation Report - Iteration 1}
\author{Aiza Ali (22i-0612) \\ Malaika Saleem (22i-0509) \\ Moimma Ali Khan (22i-1500)}
\date{October 2025}

\begin{document}

\maketitle

\tableofcontents
\clearpage

\section{Introduction}
\label{sec:introduction}

DetectifAI is an AI-powered CCTV footage analyzer designed to automate the tedious process of surveillance video analysis. This document specifies the requirements for the DetectifAI system, which serves as an intelligent surveillance assistant for security personnel, investigators, and law enforcement agencies.

The system addresses the critical problem of manual video surveillance, where security teams spend days scrolling through unedited footage to locate specific incidents. With video traffic composing over 80\% of online CCTV data worldwide, the scale of manual analysis has become impractical and error-prone.

This document is intended for multiple stakeholders:
\begin{itemize}
    \item \textbf{Developers}: Technical specifications and implementation requirements
    \item \textbf{Project Managers}: Timeline, scope, and deliverable tracking
    \item \textbf{End Users}: Security personnel, investigators, and surveillance teams
    \item \textbf{Testers}: Functional and non-functional requirement validation
    \item \textbf{Academic Evaluators}: FYP assessment and technical review
    \item \textbf{System Administrators}: Deployment and maintenance guidelines
\end{itemize}

DetectifAI automates the detection of critical security incidents including physical assault, visible weapons (guns and knives), fire detection, perimeter breaches (wall jumping), and road accidents. The system generates natural language descriptions, captures relevant keyframes, and provides searchable incident timelines to significantly reduce investigation time.

\clearpage

\section{User Classes and Characteristics}

\begin{table}[h!]
\centering
\resizebox{\textwidth}{!}{%
\begin{tabular}{p{0.2\textwidth} p{0.8\textwidth}}
\hline
\textbf{User Class} & \textbf{Description} \\ \hline

\textbf{Security Officers} & Primary users who monitor live video feeds and respond to real-time alerts. Approximately 200-500 potential users across various institutions (schools, offices, commercial buildings). These users require immediate notification of suspicious activities and quick access to incident details. They typically have basic computer skills and need an intuitive interface for rapid response scenarios. \\ \hline

\textbf{Investigation Teams} & Detectives and forensic analysts who analyze historical footage for case development and evidence collection. Expected to use the system 10-15 times per week for in-depth analysis. They require advanced search capabilities, detailed reporting features, and the ability to export evidence. These users have intermediate technical skills and need comprehensive documentation of all system interactions. \\ \hline

\textbf{Surveillance Supervisors} & Personnel responsible for overseeing multiple security operations and generating management reports. They need dashboard views of system statistics, incident trends, and team performance metrics. These users typically have administrative privileges and require both live monitoring and historical analysis capabilities. \\ \hline

\textbf{System Administrators} & Technical personnel responsible for system deployment, user management, and maintenance. They need full CRUD operations capabilities, system configuration access, and technical monitoring tools. These users have advanced technical skills and require detailed system logs and performance metrics. \\ \hline

\textbf{Emergency Responders} & External stakeholders who receive critical incident notifications and exported reports. They need simplified incident summaries with clear timestamps, locations, and event descriptions. Their interaction with the system is primarily through automated alerts and shared reports. \\ \hline

\textbf{Management Personnel} & Decision-makers who review system reports and incident summaries for policy and security planning. They require high-level dashboards, trend analysis, and executive reporting features. These users have limited direct system interaction but need comprehensive yet concise information presentation. \\ \hline

\end{tabular}%
}
\end{table}

\clearpage

\section{Functional Requirements}

\subsection{Module 1: Live Video Input \& Preprocessing}
\textbf{FR1.1} The system shall accept multiple video input formats including MP4, AVI, MOV, and real-time RTSP streams from IP cameras.

\textbf{FR1.2} The system shall automatically segment long-duration videos (>30 minutes) into manageable 10-minute chunks for efficient processing.

\textbf{FR1.3} The system shall extract keyframes at configurable intervals (default: 1 frame per second) with timestamp preservation.

\textbf{FR1.4} The system shall apply adaptive enhancement techniques (CLAHE, denoising) to improve video quality for analysis.

\textbf{FR1.5} The system shall maintain frame-to-timestamp synchronization with accuracy within 100ms.

\subsection{Module 2: Object Detection}
\textbf{FR2.1} The system shall detect fire incidents using YOLOv11 fire detection model with minimum 70\% confidence threshold.

\textbf{FR2.2} The system shall identify weapons (guns, knives) using YOLOv11 weapon detection model with minimum 75\% confidence threshold.

\textbf{FR2.3} The system shall detect human subjects and track their movements across consecutive frames.

\textbf{FR2.4} The system shall annotate detected objects with bounding boxes and confidence scores on keyframes.

\textbf{FR2.5} The system shall generate object detection events with precise timestamps and spatial coordinates.

\subsection{Module 3: Vision-Language Behavior Analysis \& Captioning}
\textbf{FR3.1} The system shall analyze keyframes to identify suspicious behaviors including physical assault, weapon handling, and unauthorized access.

\textbf{FR3.2} The system shall generate natural language descriptions for detected events with contextual details.

\textbf{FR3.3} The system shall classify events by severity level (Critical, High, Medium, Low) based on threat assessment.

\textbf{FR3.4} The system shall include person descriptions (clothing, gender, approximate age) in event captions when persons are involved.

\subsection{Module 4: Event Aggregation \& De-duplication}
\textbf{FR4.1} The system shall consolidate consecutive similar events occurring within 30-second intervals into single canonical events.

\textbf{FR4.2} The system shall eliminate duplicate detections using similarity-based deduplication with 80\% similarity threshold.

\textbf{FR4.3} The system shall prioritize events based on threat level and investigation importance.

\textbf{FR4.4} The system shall maintain event relationships and cross-references for comprehensive incident tracking.

\subsection{Module 5: Natural Language-Based Clip Retrieval}
\textbf{FR5.1} The system shall accept natural language queries (e.g., "show fire incidents last night") and return matching video segments.

\textbf{FR5.2} The system shall provide timestamp-based filtering options for refined search results.

\textbf{FR5.3} The system shall rank search results by relevance score and temporal proximity.

\textbf{FR5.4} The system shall support complex queries combining multiple criteria (time, event type, location).

\subsection{Module 6: Facial Recognition \& Image Search}
\textbf{FR6.1} The system shall capture and store facial features from persons involved in suspicious activities.

\textbf{FR6.2} The system shall detect re-occurrence of previously flagged individuals across different incidents.

\textbf{FR6.3} The system shall support image-based search allowing users to upload photos for person identification.

\textbf{FR6.4} The system shall maintain privacy-compliant facial recognition database with configurable retention policies.

\subsection{Module 7: Real-Time Monitoring Dashboard}
\textbf{FR7.1} The system shall provide live video streaming with real-time event overlay and alert notifications.

\textbf{FR7.2} The system shall generate comprehensive incident reports including keyframes, timestamps, and AI-generated descriptions.

\textbf{FR7.3} The system shall export reports in multiple formats (PDF, HTML, JSON) for external sharing.

\textbf{FR7.4} The system shall provide incident statistics dashboard with trend analysis and performance metrics.

\subsection{Module 8: User Management}
\textbf{FR8.1} The system shall support role-based access control with different permission levels (Viewer, Operator, Administrator).

\textbf{FR8.2} The system shall allow users to customize alert preferences and notification settings.

\textbf{FR8.3} The system shall maintain comprehensive audit logs of all user activities and system interactions.

\textbf{FR8.4} The system shall support secure session management with automatic timeout and multi-factor authentication.

\subsection{Module 9: Admin Module (CRUD Operations)}
\textbf{FR9.1} The system shall allow administrators to manage user accounts, camera feeds, and system configurations.

\textbf{FR9.2} The system shall provide bulk operations for user management and system maintenance.

\textbf{FR9.3} The system shall support backup and restore operations for critical system data.

\subsection{Module 10: Payment Module}
\textbf{FR10.1} The system shall integrate with Stripe payment gateway for subscription management.

\textbf{FR10.2} The system shall support multiple subscription tiers with different feature access levels.

\textbf{FR10.3} The system shall provide billing history and invoice generation capabilities.

\clearpage

\section{Non-Functional Requirements}

\subsection{Performance Requirements}
\textbf{NFR1.1} The system shall process video streams with maximum latency of 30 seconds for real-time analysis.

\textbf{NFR1.2} The system shall generate real-time alerts within 3 seconds of suspicious event detection.

\textbf{NFR1.3} The system shall support concurrent processing of up to 10 video streams simultaneously.

\textbf{NFR1.4} The system shall handle video files up to 2GB in size without performance degradation.

\subsection{Reliability Requirements}
\textbf{NFR2.1} The system shall maintain operational uptime of at least 95\% during business hours.

\textbf{NFR2.2} The system shall automatically recover from unexpected shutdowns and resume processing from the last valid checkpoint.

\textbf{NFR2.3} The system shall implement redundant storage with automatic backup every 24 hours.

\subsection{Security Requirements}
\textbf{NFR3.1} The system shall encrypt all stored video data and facial recognition information using AES-256 encryption.

\textbf{NFR3.2} The system shall authenticate all users and maintain detailed access logs for security auditing.

\textbf{NFR3.3} The system shall comply with privacy regulations regarding facial recognition and surveillance data storage.

\subsection{Scalability Requirements}
\textbf{NFR4.1} The system shall scale horizontally to support additional camera feeds and concurrent users.

\textbf{NFR4.2} The system shall optimize GPU utilization for object detection processing when available.

\subsection{Usability Requirements}
\textbf{NFR5.1} The system shall provide intuitive user interface accessible by users with basic computer skills.

\textbf{NFR5.2} The system shall respond to user interactions within 2 seconds for standard operations.

\textbf{NFR5.3} The system shall provide comprehensive help documentation and user guides.

\clearpage

\section{System Architecture Guidance}

\subsection{Architecture Style Selection}

Based on DetectifAI's requirements, the system follows a \textbf{Multi-tiered Client-Server Architecture} with \textbf{Microservices Pattern} for the backend processing modules. This hybrid approach is chosen because:

\begin{itemize}
    \item \textbf{Real-time Processing}: Event-driven architecture handles live video streams
    \item \textbf{Scalability}: Microservices allow independent scaling of AI processing modules
    \item \textbf{User Interface}: Web-based client-server for dashboard and user management
    \item \textbf{Data Flow}: Pipeline architecture for video processing workflow
\end{itemize}

\subsection{Recommended Diagrams}

\subsubsection{High-Level Architecture Diagram}
Create a \textbf{Box and Line Diagram} showing:
\begin{itemize}
    \item Frontend (React Web Application)
    \item API Gateway (Flask Backend)
    \item Video Processing Pipeline
    \item AI Detection Services (Object Detection, Facial Recognition)
    \item Event Processing Engine
    \item Database Layer (PostgreSQL + File Storage)
    \item External Services (Stripe, Notifications)
\end{itemize}

\subsubsection{Detailed Architecture Mapping}
Create a \textbf{Multi-tiered Architecture Diagram} with:
\begin{itemize}
    \item \textbf{Presentation Tier}: Web Dashboard, Mobile App
    \item \textbf{Application Tier}: Flask API, Authentication, Business Logic
    \item \textbf{Processing Tier}: Video Processing, AI Models, Event Aggregation
    \item \textbf{Data Tier}: Database, File Storage, Model Storage
\end{itemize}

\subsection{Required Design Models}

Since DetectifAI uses an \textbf{Object-Oriented Development Approach} with real-time processing, create these models:

\subsubsection{Primary Models (Required)}
\begin{enumerate}
    \item \textbf{Use Case Diagram}: Show interactions between different user classes and system functionalities
    \item \textbf{Activity Diagram}: Model the video processing workflow from input to event generation
    \item \textbf{Class Diagram}: Define core classes (VideoProcessor, EventDetector, User, etc.)
    \item \textbf{State Transition Diagram}: Model video processing states and event handling lifecycle
\end{enumerate}

\subsubsection{Secondary Models (Recommended)}
\begin{enumerate}
    \item \textbf{Sequence Diagram}: Show interaction flow for critical scenarios (real-time alert, video upload)
    \item \textbf{Component Diagram}: Detail the microservices architecture and module interactions
    \item \textbf{Deployment Diagram}: Show system deployment across different environments
\end{enumerate}

\subsection{Data Flow Modeling}

For the video processing pipeline, create \textbf{Data Flow Diagrams} at multiple levels:

\textbf{Level 0 (Context Diagram)}:
\begin{itemize}
    \item External entities: Users, Camera Systems, External APIs
    \item Single process: DetectifAI System
    \item Data flows: Video Input, Alerts, Reports, User Commands
\end{itemize}

\textbf{Level 1 (System Overview)}:
\begin{itemize}
    \item Processes: Video Input, AI Processing, Event Management, User Interface
    \item Data stores: Video Storage, Event Database, User Database
    \item Data flows between processes
\end{itemize}

\textbf{Level 2 (Detailed Processes)}:
\begin{itemize}
    \item Break down AI Processing into: Object Detection, Behavior Analysis, Event Aggregation
    \item Detail data transformations and intermediate storage
\end{itemize}

\clearpage

\section{Implementation Progress - Iteration 1}

\subsection{Completed Components}

\subsubsection{Research and Planning}
\begin{itemize}
    \item Comprehensive literature review of existing surveillance systems
    \item Technology stack selection and evaluation
    \item System architecture design and module decomposition
    \item Requirements gathering and stakeholder analysis
\end{itemize}

\subsubsection{Video Input and Preprocessing}
\begin{itemize}
    \item Implemented adaptive frame enhancement using CLAHE
    \item Developed keyframe extraction with configurable FPS
    \item Added video segmentation for large file handling
    \item Integrated motion detection and quality assessment
\end{itemize}

\subsubsection{Object Detection (Partial)}
\begin{itemize}
    \item Integrated YOLOv11 models for fire and weapon detection
    \item Implemented GPU acceleration support
    \item Developed object detection event generation
    \item Created detection confidence scoring and annotation
\end{itemize}

\subsubsection{Event Aggregation and Deduplication (Partial)}
\begin{itemize}
    \item Implemented similarity-based event deduplication
    \item Developed temporal event grouping algorithms
    \item Created threat level assessment framework
    \item Integrated event prioritization scoring
\end{itemize}

\subsubsection{User Module (Web Application)}
\begin{itemize}
    \item Developed Flask backend API structure
    \item Implemented user authentication and session management
    \item Created responsive web interface using React/Next.js
    \item Integrated role-based access control
\end{itemize}

\subsection{Technical Achievements}

\subsubsection{System Integration}
\begin{itemize}
    \item Successfully integrated all completed modules
    \item Developed comprehensive testing framework
    \item Achieved 87.5\% test success rate across all components
    \item Implemented error handling and graceful fallbacks
\end{itemize}

\subsubsection{Performance Metrics}
\begin{itemize}
    \item Video processing latency: <30 seconds for 1-minute footage
    \item Object detection accuracy: >70\% confidence on test datasets
    \item System uptime: 95\% during development testing
    \item Memory usage optimization: <2GB for standard operations
\end{itemize}

\subsection{Pending Implementation}

\subsubsection{Remaining Object Detection Features}
\begin{itemize}
    \item Complete behavior analysis for assault detection
    \item Implement wall jumping detection algorithms
    \item Develop road accident detection capabilities
\end{itemize}

\subsubsection{Advanced Features}
\begin{itemize}
    \item Full facial recognition implementation
    \item Natural language query processing
    \item Advanced reporting and analytics
    \item Real-time streaming optimization
\end{itemize}

\clearpage

\section{Testing and Validation}

\subsection{Test Strategy}
The system testing follows a comprehensive approach covering unit testing, integration testing, and system-level validation.

\subsubsection{Completed Tests}
\begin{enumerate}
    \item \textbf{Configuration Validation}: Verified all system settings and parameters
    \item \textbf{Module Import Tests}: 100\% success rate for all DetectifAI modules
    \item \textbf{Object Detection Models}: Successfully loaded and tested YOLO models
    \item \textbf{Pipeline Integration}: Complete video processing workflow validated
    \item \textbf{Facial Recognition Framework}: Placeholder implementation tested
    \item \textbf{Placeholder Modules}: All detection module frameworks functional
\end{enumerate}

\subsubsection{Performance Benchmarks}
\begin{itemize}
    \item Model loading time: Fire model (5.2MB), Weapons model (18.1MB)
    \item Processing speed: Optimized for real-time surveillance requirements
    \item Memory efficiency: Handles large video files without degradation
    \item Error handling: Graceful fallbacks when resources unavailable
\end{itemize}

\subsection{Future Testing Plans}
\begin{itemize}
    \item User acceptance testing with security personnel
    \item Performance testing under high concurrent loads
    \item Security penetration testing for data protection
    \item Integration testing with real surveillance camera systems
\end{itemize}

\clearpage

\section{Risk Assessment and Mitigation}

\subsection{Technical Risks}
\begin{itemize}
    \item \textbf{AI Model Accuracy}: Mitigated through extensive training data and continuous model improvement
    \item \textbf{Real-time Performance}: Addressed through GPU acceleration and optimization
    \item \textbf{Scalability Concerns}: Resolved via microservices architecture and cloud deployment
\end{itemize}

\subsection{Operational Risks}
\begin{itemize}
    \item \textbf{Privacy Compliance}: Implementing privacy-preserving techniques and data encryption
    \item \textbf{System Reliability}: Ensuring redundancy and automatic recovery mechanisms
    \item \textbf{User Adoption}: Providing comprehensive training and intuitive interfaces
\end{itemize}

\section{Conclusion}

The DetectifAI system has achieved significant progress in Iteration 1, with core components successfully implemented and tested. The system demonstrates strong technical feasibility and addresses the critical need for automated surveillance video analysis.

Key accomplishments include successful integration of AI-powered object detection, comprehensive video preprocessing capabilities, and a scalable system architecture. The 87.5\% test success rate validates the system's reliability and readiness for further development.

The remaining iterations will focus on completing advanced AI features, optimizing real-time performance, and conducting comprehensive user testing to ensure the system meets all stakeholder requirements.

\clearpage

\bibliographystyle{plain}
\begin{thebibliography}{10}

\bibitem{test}
IEEE Standards Association, "IEEE Recommended Practice for Software Requirements Specifications," IEEE Std 830-1998, 1998.

\bibitem{surveillance_analysis}
E. Şengönül, R. Samet, Q. Abu Al-Haija, A. Alqahtani, B. Alturki, and A. A. Alsulami, "An analysis of artificial intelligence techniques in surveillance video anomaly detection: A comprehensive survey," Applied Sciences, vol. 13, no. 8, 2023.

\bibitem{cctv_investigation}
V. Dodd, "How CCTV played a vital role in tracking Sarah Everard and her killer," The Guardian, 2021.

\bibitem{iot_safe_city}
M. K. Adnan and A. Z. Khan, "Exploring challenges of IoT-enabled safe city projects," Global Social Sciences Review, vol. VIII, pp. 259–271, Mar 2023.

\bibitem{yolo_detection}
J. Redmon and A. Farhadi, "YOLOv3: An incremental improvement," arXiv preprint arXiv:1804.02767, 2018.

\bibitem{video_analytics}
S. Bhargava and S. K. Singh, "Video surveillance system: A systematic review," in Proceedings of the International Conference on Computing, Communication and Automation, 2020.

\end{thebibliography}

\end{document}